% !TEX root = ../main.tex
\section{Store Equivalence}\label{sec:equivalence}

GenericKernel's five functions are all polymorphic in \texttt{Key}.  TreeKernel and NerveKernel call the same \texttt{GenericKernel.check} at different types.  Store Equivalence is the claim that both instantiations accept or reject the same declarations.  The file \texttt{Verify/StoreEquivalence.lean}\leanlink{https://github.com/MathNetwork/cic-kernal/blob/main/Examples/Examples/Verify/StoreEquivalence.lean\#L103} makes this concrete.

\subsection{Two instantiations of the same environment}

A polymorphic \texttt{Keys} structure bundles the 15 declaration keys.  Two concrete instantiations produce the two representations:

\begin{center}\small
\begin{tabular}{@{}p{0.47\textwidth}p{0.47\textwidth}@{}}
\toprule
\textbf{TreeKernel} (\texttt{stringKeys}) & \textbf{NerveKernel} (\texttt{uint64Keys}) \\
\midrule
\begin{lstlisting}[language={},basicstyle=\tiny\ttfamily,frame=none,backgroundcolor=\color{green!3},escapeinside={(*}{*)}]
def stringKeys : Keys String where
  nat := (*\textcolor{red!70!black}{"Nat"}*)
  zero := (*\textcolor{blue!70!black}{"Nat.zero"}*)
  succ := (*\textcolor{green!50!black}{"Nat.succ"}*)
  natRec := (*\textcolor{orange!80!black}{"Nat.rec"}*)
  add := (*\textcolor{purple!70!black}{"Nat.add"}*)
  eq := (*\textcolor{teal}{"Eq"}*)
  refl := (*\textcolor{brown}{"Eq.refl"}*)
  eqRec := (*\textcolor{violet}{"Eq.rec"}*)
  addZero := (*\textcolor{cyan!60!black}{"Nat.add\_zero"}*)
  ...
  natAddComm := (*\textcolor{magenta!70!black}{"Nat.add\_comm"}*)
\end{lstlisting}
&
\begin{lstlisting}[language={},basicstyle=\tiny\ttfamily,frame=none,backgroundcolor=\color{blue!3},escapeinside={(*}{*)}]
-- keys = H(type expression)
def uint64Keys : Keys UInt64 where
  nat := (*\textcolor{red!70!black}{1147289952..}*)
  zero := (*\textcolor{blue!70!black}{1030311651..}*)
  succ := (*\textcolor{green!50!black}{3035285639..}*)
  natRec := (*\textcolor{orange!80!black}{6708530114..}*)
  add := (*\textcolor{purple!70!black}{9302105135..}*)
  eq := (*\textcolor{teal}{1279634138..}*)
  refl := (*\textcolor{brown}{1568597419..}*)
  eqRec := (*\textcolor{violet}{2619219959..}*)
  addZero := (*\textcolor{cyan!60!black}{1498207024..}*)
  ...
  natAddComm := (*\textcolor{magenta!70!black}{7740785570..}*)
\end{lstlisting}
\\
\bottomrule
\end{tabular}
\end{center}

The function \texttt{buildEnv} is polymorphic: it constructs all 15 declarations for \emph{any} key type.  The type and proof expressions are built using parameterized helpers (\texttt{mkAdd'}, \texttt{mkEq'}, etc.)\ that reference keys from the \texttt{Keys} structure rather than hardcoded strings or integers.

\begin{center}\small
\begin{tabular}{@{}p{0.47\textwidth}p{0.47\textwidth}@{}}
\toprule
\textbf{String key calls} & \textbf{UInt64 key calls} \\
\midrule
\begin{lstlisting}[language={},basicstyle=\tiny\ttfamily,frame=none,backgroundcolor=\color{green!3}]
buildEnv stringKeys
-- internally calls:
check "Nat" natType
  none .constructor env
check "Nat.zero" zeroType
  none .constructor env
...
check "Nat.add_comm" addCommType
  (some addCommProof) .definition env
\end{lstlisting}
&
\begin{lstlisting}[language={},basicstyle=\tiny\ttfamily,frame=none,backgroundcolor=\color{blue!3},escapeinside={(*}{*)}]
buildEnv uint64Keys
-- internally calls:
check (*\textcolor{red!70!black}{1147289952..}*) natType
  none .constructor env
check (*\textcolor{blue!70!black}{1030311651..}*) zeroType
  none .constructor env
...
check (*\textcolor{magenta!70!black}{7740785570..}*) addCommType
  (some addCommProof) .definition env
\end{lstlisting}
\\
\bottomrule
\end{tabular}
\end{center}

The type expressions (\texttt{natType}, \texttt{addCommType}), proof terms (\texttt{addCommProof}), and declaration kinds (\texttt{.constructor}, \texttt{.definition}) are \emph{identical} in both columns.  Only the key values differ.

The difference emerges \emph{after} type-checking.  NerveKernel's \texttt{check} wrapper calls \texttt{GenericKernel.check} for verification, then additionally calls \texttt{shatter} to decompose the expression into nerves:

\begin{center}\small
\begin{tabular}{@{}p{0.47\textwidth}p{0.47\textwidth}@{}}
\toprule
\textbf{TreeKernel: check only} & \textbf{NerveKernel: check + shatter} \\
\midrule
\begin{lstlisting}[language={},basicstyle=\tiny\ttfamily,frame=none,backgroundcolor=\color{green!3}]
-- 1. type-check
let env := GenericKernel.check
  "Nat.add" addType (some addValue)
  .definition env

-- done. declaration stored in
-- HashMap String Decl
-- no further processing
\end{lstlisting}
&
\begin{lstlisting}[language={},basicstyle=\tiny\ttfamily,frame=none,backgroundcolor=\color{blue!3},escapeinside={(*}{*)}]
-- 1. type-check (same logic)
let env := GenericKernel.check
  (*\textcolor{purple!70!black}{9302105135..}*) addType (some addValue)
  .definition env

-- 2. shatter: decompose into nerves
let (_, net) := shatter addValue net
-- 21 tree nodes -> 15 nerves
-- Nat stored once (was 4 copies)
-- result: HashMap UInt64 AstroNerve
\end{lstlisting}
\\
\bottomrule
\end{tabular}
\end{center}

TreeKernel stores declarations and stops.  NerveKernel additionally decomposes each expression into content-addressed nerves, producing the hypergraph structure described in Section~\ref{sec:tree-to-nerve}.  The type-checking step is identical; the shatter step is what creates the nerve store.

\subsection{Verification via native\_decide}

Three theorems confirm that both instantiations produce the same result:

\begin{lstlisting}[language={},morekeywords={theorem,by}]
-- TreeKernel succeeds on all 15 declarations
theorem lean_build_succeeds :
    (buildEnv stringKeys).isSome = true := by native_decide

-- NerveKernel succeeds on all 15 declarations
theorem astro_build_succeeds :
    (buildEnv uint64Keys).isSome = true := by native_decide

-- Both produce environments of equal size
theorem lean_astro_same_size :
    (buildEnv stringKeys).get!.size
    = (buildEnv uint64Keys).get!.size := by native_decide
\end{lstlisting}

Each theorem is proved by \texttt{native\_decide}, which compiles the expression to native code and evaluates it at compile time.  This is concrete verification on a specific 15-declaration environment.
